% Preámbulo común a las dos ediciones (Kobo y A4).
% Cada edición define antes: clase, geometry, \edicion y \figescala.

\usepackage[sc]{mathpazo}
\usepackage{fontspec}
\setmainfont{TeX Gyre Pagella}
\setsansfont{TeX Gyre Heros}[Scale=0.92]
\setmonofont{DejaVu Sans Mono}[Scale=0.82]
\usepackage{babel}
\babelprovide[import,main]{spanish}
\usepackage{amsmath,amssymb,mathtools}
\usepackage{microtype}
\usepackage[dvipsnames]{xcolor}
\usepackage{graphicx}
\usepackage{booktabs,tabularx,array,multirow}
\usepackage{enumitem}
\usepackage[most]{tcolorbox}
\usetikzlibrary{positioning,arrows.meta}
\usepackage{titlesec}
\usepackage{fancyhdr}
\usepackage{caption}
\usepackage{float}
\usepackage{csquotes}
\usepackage{bookmark}
\usepackage{hyperref}
\hypersetup{colorlinks=true,linkcolor=azul,urlcolor=azul,citecolor=azul,
  pdftitle={Derivados financieros desde cero: el caso PETROPERÚ},
  pdfauthor={Libro de estudio},bookmarksnumbered=true}

% ---------- Colores (contrastados para pantallas Kaleido) ----------
\definecolor{azul}{HTML}{1D4E9E}
\definecolor{rojo}{HTML}{C0391B}
\definecolor{verde}{HTML}{1F7A38}
\definecolor{oro}{HTML}{A86F00}
\definecolor{morado}{HTML}{6E2A8E}
\definecolor{gris}{HTML}{4A4A4A}

% ---------- Párrafos y listas ----------
\setlength{\parindent}{0pt}
\setlength{\parskip}{0.45em plus 0.15em}
\linespread{1.06}
\setlist{nosep,leftmargin=1.1em,topsep=2pt,itemsep=1pt}
\setlength{\emergencystretch}{3em}
\tolerance=2000
\raggedbottom
\setlength{\tabcolsep}{3pt}
\captionsetup{font=small,labelfont={bf,color=azul},skip=3pt,justification=raggedright,singlelinecheck=false}
\setcounter{tocdepth}{1}
\setcounter{secnumdepth}{2}

% ---------- Títulos ----------
\titleformat{\part}[display]{\centering\sffamily\bfseries\color{azul}}{\Large Parte \thepart}{6pt}{\LARGE}
\titleformat{\chapter}[display]{\sffamily\bfseries\color{azul}}{\small\MakeUppercase{Capítulo} \thechapter}{2pt}{\Large\raggedright}[\vspace{-2pt}{\color{azul}\titlerule[0.8pt]}]
\titlespacing*{\chapter}{0pt}{0pt}{8pt}
\titleformat{\section}{\sffamily\bfseries\large\color{azul}\raggedright}{\thesection}{0.5em}{}
\titlespacing*{\section}{0pt}{10pt}{3pt}
\titleformat{\subsection}{\sffamily\bfseries\normalsize\color{gris}\raggedright}{\thesubsection}{0.5em}{}
\titlespacing*{\subsection}{0pt}{8pt}{2pt}
\titleformat{\paragraph}[runin]{\bfseries\color{azul}}{}{0pt}{}[.]
\titlespacing*{\paragraph}{0pt}{4pt}{0.5em}

% ---------- Cabecera ----------
\pagestyle{fancy}
\fancyhf{}
\fancyhead[L]{\scriptsize\sffamily\color{gris}\nouppercase{\leftmark}}
\fancyhead[R]{\scriptsize\sffamily\color{gris}\thepage}
\renewcommand{\headrulewidth}{0.3pt}
\renewcommand{\chaptermark}[1]{\markboth{\thechapter.\ #1}{}}
\fancypagestyle{plain}{\fancyhf{}\fancyhead[R]{\scriptsize\sffamily\color{gris}\thepage}\renewcommand{\headrulewidth}{0pt}}

% ---------- Recuadros ----------
\tcbset{base/.style={enhanced,breakable,boxrule=0.6pt,arc=1.5pt,left=3pt,right=3pt,top=2pt,bottom=2pt,
  fonttitle=\sffamily\bfseries\small,before skip=6pt,after skip=6pt,parbox=false}}
\newtcolorbox{definicion}[1][]{base,colback=azul!6,colframe=azul,title={Definición\if\relax\detokenize{#1}\relax\else: #1\fi}}
\newtcolorbox{intuicion}[1][]{base,colback=verde!6,colframe=verde,title={Intuición\if\relax\detokenize{#1}\relax\else: #1\fi}}
\newtcolorbox{importante}[1][]{base,colback=oro!8,colframe=oro,title={\if\relax\detokenize{#1}\relax Concepto clave\else #1\fi}}
\newtcolorbox{cuidado}[1][]{base,colback=rojo!5,colframe=rojo,title={Error frecuente\if\relax\detokenize{#1}\relax\else: #1\fi}}
\newtcolorbox{informe}[1][]{base,colback=gris!6,colframe=gris,title={En el informe\if\relax\detokenize{#1}\relax\else: #1\fi}}
\newtcolorbox{adicional}[1][]{base,colback=morado!5,colframe=morado,title={Contenido adicional\if\relax\detokenize{#1}\relax\else: #1\fi}}
\newtcolorbox{resumen}{base,colback=azul!4,colframe=azul,title={Resumen del capítulo}}
\newtcolorbox{objetivos}{base,colback=white,colframe=azul,boxrule=0.8pt,title={Objetivos de aprendizaje}}
\newtcolorbox{prerreq}{base,colback=white,colframe=gris,title={Prerrequisitos}}

\newcounter{ejemplo}[chapter]
\renewcommand{\theejemplo}{\thechapter.\arabic{ejemplo}}
\newtcolorbox{ejemplo}[1]{base,colback=white,colframe=verde,boxrule=0.8pt,
  title={\refstepcounter{ejemplo}Ejemplo \theejemplo: #1}}

% Ficha de ecuación con 8 pasos
\newtcolorbox{ficha}[1]{base,colback=white,colframe=azul,boxrule=0.9pt,title={Ficha de la ecuación: #1}}
\newcommand{\paso}[1]{\par\smallskip{\sffamily\bfseries\small\color{azul}#1.}\ }

% Ejercicios
\newcounter{ejer}[chapter]
\renewcommand{\theejer}{\thechapter.\arabic{ejer}}
\newcommand{\ejercicio}[1][]{\par\medskip\refstepcounter{ejer}{\sffamily\bfseries\color{azul}Ejercicio \theejer}\if\relax\detokenize{#1}\relax\else\ {\small(#1)}\fi.\ }
\newcommand{\solucion}[1]{\par\medskip{\sffamily\bfseries\color{verde}Solución #1.}\ }

% Atajos
\newcommand{\US}{US\$}
\newcommand{\MM}{\,millones}
\newcommand{\figura}[3][0.98]{\begin{figure}[H]\centering\includegraphics[width=\figescala\dimexpr #1\linewidth\relax]{fig/#2}\caption{#3}\label{fig:#2}\end{figure}}
\newcommand{\pb}{\,\text{pb}}
\DeclareMathOperator{\VP}{VP}
\DeclareMathOperator{\VF}{VF}
\newcommand{\E}{\mathbb{E}}
